\begin{proof}[Proof of \cref{obj:thm:k-label-rate}]
The overlap condition gives \(\ell>0\). First fix \(H\in\mathcal H^{\mathrm{lb}}_{\varepsilon,m_f}\) and an arbitrary measurable \(g\in\mathcal G_K\). For each label \(r\), set
\[
C_r=\{e\in\mathcal E:g(e)=r\},
\qquad h_r=\operatorname{Leb}(C_r).
\]
The measurable cells partition \(\mathcal E\), including any empty cells, so \(h_r\geq0\) and \(\sum_r h_r=\ell\). % lean: realScoreCell_intervalPartition

For every \(\zeta\in\mathbb R\) and \(t\geq0\), the points within distance \(t\) of \(\zeta\) occupy an interval of length \(2t\). Thus
\[
\operatorname{Leb}\{e\in C_r:|e-\zeta|>t\}\geq h_r-2t,
\]
and the layer-cake identity gives
\[
\int_{C_r}|e-\zeta|\,de
\geq\int_0^{h_r/2}(h_r-2t)\,dt
=\frac{h_r^2}{4}.
\] % lean: boundedCell_absDeviation_integral_lower

Since \(H(de)=f(e)\,de\) with \(f\geq m_f\) almost everywhere, it follows for every \(\zeta\in\mathbb R\) that
\[
\int_{C_r}|e-\zeta|\,H(de)
\geq m_f\int_{C_r}|e-\zeta|\,de
\geq \frac{m_fh_r^2}{4}.
\] % lean: scoreCell_absDeviation_lower

Write \(p_1(e)=e\) and \(p_0(e)=1-e\), as in \cref{obj:thm:all-label-ambiguity}. For any \(c\in[1/(1-\varepsilon),1/\varepsilon]\), define the two score centers by
\[
z_{1,c}=c^{-1},\qquad z_{0,c}=1-c^{-1}.
\]
Here \(c>0\), and direct algebra gives
\[
|1-cp_a(e)|=c|e-z_{a,c}|,
\qquad a\in\{0,1\}.
\] % lean: reciprocalDeviation_eq_scoreDistance

Because \(c\geq1/(1-\varepsilon)\), the preceding cell bound yields, for either arm,
\[
\int_{C_r}|1-cp_a(e)|\,H(de)
\geq \frac{m_fh_r^2}{4(1-\varepsilon)}.
\] % lean: cellAbsoluteDeviation_lower_cellLengthSq

By the exact cell formula in \cref{obj:thm:all-label-ambiguity},
\[
D_H(g)
=\sum_r\sum_{a\in\{0,1\}}
\inf_{c\in[1/(1-\varepsilon),\,1/\varepsilon]}
\int_{C_r}|1-cp_a(e)|\,H(de)
\geq \frac{m_f}{2(1-\varepsilon)}\sum_r h_r^2.
\] % lean: worstCaseAmbiguity_eq

There are \(K\) labels, so Cauchy–Schwarz gives
\[
\sum_r h_r^2
\geq \frac{(\sum_r h_r)^2}{K}
=\frac{\ell^2}{K}.
\] % lean: sum_cellMass_sq_lower
Consequently every \(g\in\mathcal G_K\) satisfies
\[
D_H(g)\geq\frac{m_f\ell^2}{2(1-\varepsilon)K}.
\] % lean: worstCaseAmbiguity_finiteLabel_lower
The class \(\mathcal G_K\) is nonempty because \(K\geq1\) permits a constant measurable release. Taking the infimum as in \cref{obj:def:optimal-k-ambiguity} proves the lower bound for \(\Delta_K(H)\). % lean: optimalKAmbiguity_finiteLabel_lower

Now let \(H\) be any probability law on \(\mathcal E\), and take the equal-width release \(g\) in the statement. It is measurable: the score transformation inside the floor is continuous, the floor is measurable, and clipping its value to the finite label set preserves measurability. % lean: equalWidthRelease_measurable

Index its labels by \(r=1,\ldots,K\), write \(C_r=\{e\in\mathcal E:g(e)=r\}\), and define
\[
z_r=\varepsilon+\left(r-\frac12\right)\frac{\ell}{K}.
\]
For \(e\in C_r\), the floor rule places \(e\) between \(\varepsilon+(r-1)\ell/K\) and \(\varepsilon+r\ell/K\). At the right endpoint of \(\mathcal E\), clipping assigns the last label, for which the same closed bounds hold. Hence \(z_r\in\mathcal E\) and
\[
|e-z_r|\leq\frac{\ell}{2K}\qquad(e\in C_r).
\] % lean: equalWidthRelease_cell_subset_closedBin

The candidates \(z_r^{-1}\) and \((1-z_r)^{-1}\) lie in \([1/(1-\varepsilon),1/\varepsilon]\). Moreover,
\[
z_r(1-z_r)-\varepsilon(1-\varepsilon)
=(z_r-\varepsilon)(1-\varepsilon-z_r)\geq0.
\]
Thus, for \(e\in C_r\),
\[
\left|1-\frac e{z_r}\right|
+\left|1-\frac{1-e}{1-z_r}\right|
=\frac{|e-z_r|}{z_r(1-z_r)}
\leq \frac{|e-z_r|}{\varepsilon(1-\varepsilon)}.
\] % lean: sum_arm_reciprocalDeviation_le

Using these candidates in the two cell infima, then applying the midpoint bound, gives
\[
\sum_{a\in\{0,1\}}
\inf_{c\in[1/(1-\varepsilon),\,1/\varepsilon]}
\int_{C_r}|1-cp_a(e)|\,H(de)
\leq
\frac{1}{\varepsilon(1-\varepsilon)}
\int_{C_r}|e-z_r|\,H(de)
\leq
\frac{\ell}{2\varepsilon(1-\varepsilon)K}H(C_r).
\] % lean: sum_cellAbsoluteDeviation_le_centeredScoreIntegral

The cells partition \(\mathcal E\) and \(H(\mathcal E)=1\). Summing the cell bounds therefore yields
\[
D_H(g)\leq
\frac{\ell}{2\varepsilon(1-\varepsilon)K}
\sum_r H(C_r)
=
\frac{\ell}{2\varepsilon(1-\varepsilon)K}.
\] % lean: equalWidthRelease_worstCaseAmbiguity_le_all

Finally, the exact cell formula shows \(D_H(g')\geq0\) for every measurable release \(g'\), while the measurable \(g\) above belongs to \(\mathcal G_K\). Hence the defining infimum satisfies
\[
\Delta_K(H)\leq D_H(g)
\leq\frac{\ell}{2\varepsilon(1-\varepsilon)K}.
\] % lean: optimalKAmbiguity_bounds
\end{proof}
